\head{第二十版整合范围与核查方法}
\audit{V20-INT01}{本次实际处理范围}{第一篇第12节完整法英回校；其他已译内容按文件继承}
基线是公开仓库提交\texttt{997597a5476cb07338a9b2feef61077ae795774e}中的\texttt{current/EGA\_v19/}。本版从第五卷预备笔记第一篇第12节标题开始，逐段处理法稿物理64--73页，对读英编稿43页中段至48页末；法稿67和68页是同一原页的重复扫描，两个扫描都检查，正文不重复排印。法稿74页和英编稿49页用于验证第13节边界，不记入本轮回校范围。

四个新正文模块按“关联概形与相对Grassmann描述”“既有命题向高余维截面的推广”“奇异关联概形与退化层”“超奇异点、线性截面族和一般中心”顺序相接。旧\texttt{v\_dot\_middle.tex}完整保存，退出当前输入链；其中第13节的全部字节另存为继承模块并继续输入，随后仍接旧第14、15节。没有把第13节及以后已有译文丢掉，也没有把它们算作本轮已经重新审定。

Dot117--119的独立工程原ZIP位于公开的\texttt{archive\_\allowbreak supplement\_\allowbreak20261125/\allowbreak dot/}。本轮实际解包其候选译文、核查与来源资料，但采用原页对读后的结论；其错误归属、过强解释或引用错误不会因来源于独立批稿就无条件沿用。Dot109--116此前对应的回校范围继续由第十九版源码承接，不重复导入。Dot69--108的40个模块各保留一次输入。

本轮性质是已有稿件的回校、缺漏恢复、订正和明确标注的译者补证，并非再从头新增整节译文。成品PDF页数包括整个历史合订稿、目录和脚注，不能据此换算本轮新译原文数量。

\audit{V20-INT02}{错误归属与证明身份}{原作者、英文转写、旧译和新补证分别标明}
法稿本来正确、英文转写失误的内容，据原页恢复并纠正旧报告的归责。法稿本身的矛盾或数学错误，正文用有依据的表述，原说法和反例入脚注。字形仍有歧义的手改符号与数学上可确定的正确界限分别说明，不把数学推断冒充已辨明的字迹。法稿删划、编写建议、英译编辑注和历史译者补证各按其原身份保存。

本轮最明确的新反例涉及第12节末：二次Veronese圆锥曲线的一般直线束有两个坏参数，否定原稿对整个坏参数集几何不可约性的断言；它不否定该例中余维1和每个坏纤维唯一普通奇点的性质。末尾一般中心的余维界另给完整旗概形维数补证，但不由此推得交集几何不可约。详细公式和适用条件见本报告V20-FAM-E02、E03。

其余重点是相对Grassmann维数中的秩约定、正则列与Koszul正则性的适用条件、“自8.6起”的方向性恢复，以及奇异零点退化结构与单个纤维非光滑概形的区别。报告逐项区分已经证明、仅作反例排除和依赖原稿未完成推广的内容。

\audit{V20-INT03}{工程核验及其边界}{源码、编译、逐页一致性与视觉抽查分别记录}
本版沿用真正的\texttt{ctexbook}类及其默认字体、字号和正文版心，所有公式与交换图保持TeX排版。源文件继承核验、当前输入链及单文件展开检查、四份XeLaTeX构建日志、模块版与独立单文件版的逐页文字和像素比对、最终变动页的实际渲染抽查分别输出记录。具体执行结果、页数和成品散列以本版\texttt{checks/}、\texttt{PROGRESS.md}和\texttt{ARTIFACTS.json}为准；未实际执行的检查不得标成通过。

历史I、II、III章、第0章以及本轮范围以外的第五卷稿件没有重新逐句审定。正式第IV章主体仍未连续译入；英文预备稿到达物理112页末也不能证明法文171页全部逐页审结。全部书籍尚未完成。下一默认回校位置是法稿物理74页、英编稿49页的第13节标题及定义13.1。

历史原始归档维持原状。此前档案核查已经明列7个早期独立工程原ZIP无法恢复，不能因制作新工程或发布新版本而删除这一事实。当前项目包与包含底本、Dot和可恢复历史的全库工程包继续明确区分。
